\subsection{From database IDs to list-ordered positions}

List ordering makes the vectors assigned to one IVF centroid contiguous.
An input row has identifier $i\in\{0,\ldots,N-1\}$; coarse assignment
uses computed centroid scores with the mathematical target
\begin{equation}
 a(i)=\operatorname*{arg\,min}_{0\leq\ell<n_{\rm list}}
       \|\Pi x_i-\mu_\ell\|_2^2.
 \label{eq:4-assignment}
\end{equation}
The implementation uses the numerical assignment path described above.
Primary and residual encoding do not subtract $\mu_{a(i)}$: list
membership restricts the search universe independently of those
reconstructions.

The builder sorts assignments together with input IDs. Let $\rho(p)$
be the input ID at sorted physical position $p$, and $o_\ell$ the start
of list $\ell$. The completed arrays satisfy
\begin{equation}
 0=o_0\leq o_1\leq\cdots\leq o_{n_{\rm list}}=N,
 \qquad p\in[o_\ell,o_{\ell+1})\Longrightarrow a(\rho(p))=\ell.
 \label{eq:4-offsets}
\end{equation}
Lower-bound searches find each list's first position; equal offsets
therefore mark an empty list:
\[
 o_\ell=o_{\ell+1}\Longleftrightarrow\text{list }\ell\text{ is empty}.
\]
Positions stay fixed in a completed index but can change after
rebuilding. All arrays use the same $\rho(p)$, and IDs map
$p\mapsto\rho(p)$ back to input rows. Two-pass add temporarily builds
$\rho^{-1}$ to write new encodings directly at their destination
(scatter). Primary ties use physical positions; final ties use database
IDs.

\par\medskip\noindent\begin{minipage}{\linewidth}
\captionsetup{hypcap=false}
\centering
% Independent arrays share a physical position; the band replaces connector lines.
\begin{tikzpicture}[x=1cm,y=1cm,font=\small\linespread{1}\selectfont,
 cell/.style={draw=ThesisNavy!45,minimum width=.86cm,minimum height=.48cm,inner sep=1pt},
 panel/.style={draw=ThesisNavy!35,rounded corners=2pt,fill=ThesisNavy!2}]
\draw[panel] (0,-1.85) rectangle (15.1,-.85);
\draw[panel] (0,-4.15) rectangle (15.1,-2.15);
\draw[panel] (0,-5.55) rectangle (15.1,-4.55);
\draw[panel] (0,-6.95) rectangle (15.1,-5.95);
\fill[ThesisNavy!12] (6.22,-6.90) rectangle (7.08,.3);
\node[font=\small\bfseries] at (6.65,1.12) {$p=3$};
\node[anchor=west,align=left] at (.15,.03) {$o=[0,3,3,8]$\\list 1 is empty};
\node[anchor=west,font=\small\bfseries] at (11.2,.03) {Byte offset};
\foreach \p in {0,...,7} {\node[cell] at (3.95+.90*\p,0) {$\p$};}
\draw[decorate,decoration={brace,amplitude=3pt}] (3.52,.32)--(6.18,.32) node[midway,above=4pt] {list 0};
\draw[decorate,decoration={brace,amplitude=3pt}] (6.22,.32)--(10.68,.32) node[midway,above=4pt] {list 2};
\node[anchor=west,align=left] at (.2,-1.35) {\textbf{Input IDs}\\$\mathrm{IDs}[p]=\rho(p)$};
\foreach \p/\id in {0/5,1/1,2/7,3/0,4/6,5/3,6/4,7/2} {\node[cell] at (3.95+.90*\p,-1.35) {$\id$};}
\node[anchor=west] at (11.2,-1.35) {$4p$};
\node[anchor=west,align=left] at (.2,-2.95) {\textbf{Primary words}\\4 codes per word};
\foreach \j/\yy in {0/-2.70,1/-3.30} {
 \foreach \p in {0,...,7} {\node[cell] at (3.95+.90*\p,\yy) {$\mathsf w_{\j,\p}$};}}
\node[anchor=west] at (11.2,-2.98) {$4(jN+p)$};
\node at (7.1,-3.83) {Adjacent candidates (threads) read consecutive words};
\node[anchor=west,align=left] at (.2,-5.05) {\textbf{Residual rows}\\$b_r$ bytes per row};
\foreach \p in {0,...,7} {\node[cell] at (3.95+.90*\p,-5.05) {$r_{\p}$};}
\node[anchor=west] at (11.2,-5.05) {$pb_r+t$};
\node[anchor=west,align=left] at (.2,-6.45) {\textbf{Corrections}\\FP32};
\foreach \p in {0,...,7} {\node[cell] at (3.95+.90*\p,-6.45) {$\gamma[\p]$};}
\node[anchor=west] at (11.2,-6.45) {$4p$};
\end{tikzpicture}

\captionof{figure}[Physical GPU index layout]{Physical GPU index layout.}
\label{fig:4-layout}
\end{minipage}\par\medskip

The example uses $N=8$ and offsets $o=[0,3,3,8]$, so list~1 is empty. The shaded column marks $p=3$: four separate arrays use the same physical position, and byte offsets are relative to each array's base. At fixed word group $j$, adjacent candidates occupy consecutive 32-bit words, allowing threads scanning adjacent candidates to coalesce their reads. In the eight-bit residual configuration, each row holds the $b_r$ code bytes of one vector.

\subsection{Exact addresses and representation widths}

The primary representation groups candidates by subspace. For primary
byte $Z[m,p]$, its address is
\[
 \operatorname{offset}_{\rm byte}(Z[m,p])=mN+p.
\]
Final packing combines four subspace bytes of one vector into each word:
\begin{equation}
 \mathsf w[j,p]=\sum_{t=0}^{3}2^{8t}Z[4j+t,p],\qquad
 \operatorname{addr}_{\rm byte}(\mathsf w[j,p])=4(jN+p).
 \label{eq:4-primary-address}
\end{equation}
Here $j=0,\ldots,M/4-1$. $\text{word-layout build}\Longrightarrow M\bmod4=0$;
other $M$ are rejected, while a separate byte-layout compilation retains
the byte representation. Each array occupies $NM$ bytes; repacking
temporarily holds both copies, costing $2NM$. The one-pass builder
reorders existing codes (gather) into a vector-major $[N,M]$ array,
transposes it and repacks, freeing obsolete copies. Two-pass construction
scatters directly into subspace-major storage before the same repack.
The search-time lane mapping is discussed in \Cref{sec:5-3}.

Residual codes remain row-major, with $b_r=\lceil dB_r/8\rceil$ bytes per
vector. In the eight-bit configuration, coordinate $t$ of row $p$ is
at byte offset $pb_r+t$. In the four-bit path it is in byte
$pb_r+\lfloor t/2\rfloor$, with even coordinates in the low four bits
(low nibble) and odd coordinates in the high four bits (high nibble).
A four-byte correction per physical row stores
\begin{equation}
 \gamma[ p ]=2\langle\widehat p(x_{\rho(p)}),
                              \widehat r(x_{\rho(p)})\rangle.
 \label{eq:4-correction}
\end{equation}
The correction must move with its row's codes during gathering or
scattering. List-order scans can then fetch residuals only for retained
candidates. Refinement uses compressed reconstructions; raw FP32
reranking would require an additional $4Nd$ device bytes. Chapter 7
evaluates the resulting recall.
